% BHU PYQ -- Commerce: Financial Accounting (2025-26 NEP)
\documentclass[11pt,a4paper]{article}
\usepackage[a4paper,top=2.5cm,bottom=2.5cm,left=2.8cm,right=2.8cm,headheight=14pt]{geometry}
\usepackage{amsmath,amssymb}


\usepackage{lmodern,microtype,fancyhdr,enumitem,hyperref,lastpage,parskip,booktabs}

\hypersetup{colorlinks=true,urlcolor=black,linkcolor=black,pdftitle={COMMJ11: Financial Accounting -- BHU NEP 2025-26}}

\pagestyle{fancy}\fancyhf{}
\renewcommand{\headrulewidth}{0.4pt}\renewcommand{\footrulewidth}{0.4pt}
\fancyhead[L]{\small \textit{Banaras Hindu University}}
\fancyhead[R]{\small \textit{COMMJ11: Financial Accounting (2025-26)}}
\fancyfoot[C]{\small Page \thepage\ of \pageref{LastPage}}

\newlist{parts}{enumerate}{2}
\setlist[parts,1]{label=(\alph*),leftmargin=*,itemsep=4pt,topsep=3pt}
\setlist[parts,2]{label=(\roman*),leftmargin=*,itemsep=2pt,topsep=2pt}

\newcommand{\pts}[1]{\hfill \textit{[#1]}}
\newcommand{\rupee}{Rs.~}

\begin{document}

\begin{center}
  {\large\bfseries BANARAS HINDU UNIVERSITY}\\[2pt]
  {\normalsize B. Com. (Hons.) Semester I Examination, 2025-26 (NEP)}\\[6pt]
  \rule{0.8\linewidth}{0.5pt}\\[6pt]
  {\large\bfseries COMMERCE}\\[4pt]
  {\large\bfseries Paper Code: COMMJ11 : Financial Accounting}\\[4pt]
  {\normalsize Time: 2:30 Hours \hfill Full Marks: 70}\\[2pt]
  {\small REF NO. CECONF/2025-26/PS/24}
\end{center}
\rule{\linewidth}{0.4pt}
\smallskip

\noindent \textit{Instructions: Answer \textbf{all} questions. All questions carry equal marks (14 marks each). Write your Roll No.\ at the top immediately on receipt of this question paper.}

\smallskip
\noindent \textit{सभी प्रश्नों के उत्तर दीजिए। सभी प्रश्नों के अंक समान हैं।}

\bigskip

\begin{enumerate}[label=\textbf{\arabic*.},leftmargin=*,itemsep=14pt]

  \item 
  \begin{parts}
    \item Explain in brief Provisions and Reserves.\pts{7}
    
    संक्षेप में प्रावधान तथा संचय को समझाइए।

    \item Explain Trial Balance. State the errors which are not disclosed by Trial Balance.\pts{7}
    
    तलपट की व्याख्या कीजिए। उन अशुद्धियों को बताइए जो तलपट द्वारा प्रकट नहीं होती हैं।
  \end{parts}

  \begin{center}\textbf{OR / अथवा}\end{center}

  Prepare Bank Reconciliation Statement from the following particulars and show the balance as per pass book:\pts{14}

  निम्नलिखित विवरणों से बैंक समाधान विवरण तैयार कीजिए तथा पास बुक के अनुसार शेष दर्शाए:

  \begin{parts}
    \item Cash book of Singh Traders showed an overdraft of \rupee~82,000.
    \item Cheque issued but not encashed \rupee~23,000.
    \item Cheque paid into Bank but not credited \rupee~30,000.
    \item Interest charged on overdraft \rupee~1,200.
    \item Bank Charges \rupee~500.
    \item Interest on debentures collected by Bank \rupee~6,000.
    \item Insurance Premium paid by Bank \rupee~2,200.
    \item A customer paid to Singh Traders directly into Bank \rupee~10,000.
  \end{parts}

  \item Prepare Trading and Profit and Loss Account for the year ended 31st December, 2024 and Balance Sheet as on that date from the following Trial Balance:\pts{14}

  \begin{center}
    \small
    \begin{tabular}{lr | lr}
      \toprule
      \textbf{Debit Particulars} & \textbf{Amount (\rupee)} & \textbf{Credit Particulars} & \textbf{Amount (\rupee)} \\
      \midrule
      Land \& Building & 70,000 & Sales & 2,00,000 \\
      Plant \& Machinery & 60,000 & Purchase Return & 12,000 \\
      Loose tools & 10,000 & Capital & 2,50,000 \\
      Bills Receivable & 15,000 & Creditors & 95,000 \\
      Opening Stock & 50,000 & Bank Loan & 16,000 \\
      Purchases & 1,40,000 & Commission & 2,000 \\
      Wages & 40,000 & & \\
      Carriage & 5,000 & & \\
      Salaries & 25,000 & & \\
      Rents \& Rates & 5,000 & & \\
      Discount & 4,000 & & \\
      Cash & 8,000 & & \\
      Debtors & 80,000 & & \\
      Bad Debts & 4,000 & & \\
      Furniture & 20,000 & & \\
      Advertising & 8,000 & & \\
      Sales Return & 15,000 & & \\
      Goodwill & 16,000 & & \\
      \midrule
      \textbf{Total} & \textbf{5,75,000} & \textbf{Total} & \textbf{5,75,000} \\
      \bottomrule
    \end{tabular}
  \end{center}

  \textbf{Adjustments:}
  \begin{parts}
    \item Closing Stock \rupee~78,000.
    \item Prepaid Advertising \rupee~1,500.
    \item Interest on Bank Loan @ 10\% p.a.
    \item Depreciate Plant \& Machinery @ 10\% and Loose tools @ 15\% p.a.
    \item Provide for doubtful debts @ 5\% p.a.
    \item Free samples distributed \rupee~800.
  \end{parts}

  \begin{center}\textbf{OR / अथवा}\end{center}

  \begin{parts}
    \item Explain Bills of Exchange and its various parties. What is the difference between Accommodation Bill and Promissory Note?\pts{7}
    \item Give the journal entries in the books of Drawer and Drawee.\pts{7}
  \end{parts}

  \item $X$ and $Y$ share profits in the proportion of $5/8$ and $3/8$. The Balance Sheet of the firm as on 31st March, 2025 was as under:\pts{14}

  \begin{center}
    \small
    \begin{tabular}{lr | lr}
      \toprule
      \textbf{Liabilities} & \textbf{Amount (\rupee)} & \textbf{Assets} & \textbf{Amount (\rupee)} \\
      \midrule
      Creditors & 25,000 & Cash at Bank & 11,200 \\
      General Reserve & 20,000 & Bills Receivable & 12,800 \\
      Capital Accounts: & & Debtors & 20,000 \\
      ~~$X$: 75,000 & & Stock & 35,000 \\
      ~~$Y$: 60,000 & 1,35,000 & Furniture & 21,000 \\
      & & Machinery & 30,000 \\
      & & Building & 50,000 \\
      \midrule
      \textbf{Total} & \textbf{1,80,000} & \textbf{Total} & \textbf{1,80,000} \\
      \bottomrule
    \end{tabular}
  \end{center}

  On the above date $Z$ was admitted on the following terms:
  \begin{parts}
    \item $Z$ was to get $1/5$ share in profits and he would pay \rupee~50,000 as capital and \rupee~16,000 as share of Goodwill.
    \item Machinery was to be depreciated by 10\% and building was to be appreciated by 20\%. Stock was valued at 25\% above cost. Stock to be shown in new firm at cost price.
    \item Liability for repair to furniture \rupee~600.
    \item Capital Accounts of old partners were to be adjusted in new profit sharing ratio.
  \end{parts">
  Prepare Revaluation Account, Capital Accounts and initial Balance Sheet of the new firm.

  \begin{center}\textbf{OR / अथवा}\end{center}

  \begin{parts}
    \item What is meant by Revaluation Account? Differentiate between Revaluation Account and Profit \& Loss Account.\pts{7}
    \item Explain the different methods of treating premium paid on Joint Life Policies. Show the entries also.\pts{7}
  \end{parts}

  \item RK Ltd., registered with a nominal capital of 10,000 shares of \rupee~100 each, issued 5,000 shares at a premium of \rupee~10 per share, payable as \rupee~20 on Application, \rupee~50 on Allotment, \rupee~20 on First Call and \rupee~20 on Second \& Final Call. Applications were received for 6,000 shares and pro-rata allotment was made. One shareholder having 300 shares could not pay allotment money and on his subsequent failure to pay first call, his shares were forfeited. Another shareholder having 200 shares failed to pay two calls and his shares were accordingly forfeited. These forfeited shares were reissued for \rupee~90 per share as fully paid up.\pts{14}

  Pass the necessary journal entries of the above transactions.

  \begin{center}\textbf{OR / अथवा}\end{center}

  Give the proforma of Profit and Loss Account and Balance Sheet as per Indian Companies Act with imaginary figures.\pts{14}

  \item Answer the following in not more than 75 words each:\pts{4 $\times$ 3.5 = 14}
  \begin{parts}
    \item Accounting Concepts (लेखांकन अवधारणाएँ)
    \item Capital Expenditure (पूँजीगत व्यय)
    \item Adjustments (समायोजन)
    \item Issue and redemption of debentures (ऋणपत्रों का निर्गमन तथा शोधन)
  \end{parts}

\end{enumerate}

\end{document}
